\section{Experiments}
\label{sec:Experiments}

We conduct extensive experiments on a real-world robotic platform to evaluate the performance of our proposed method, GenSplat. Our evaluation is structured around four central research questions (RQs) designed to assess our contributions systematically:

\begin{enumerate}[label=RQ\arabic*:, leftmargin=*]
    \item How effectively does expanding the observational manifold via GenSplat improve policy robustness to out-of-distribution camera perturbations?
    % How effectively does GenSplat augmentation improve policy robustness to camera perturbations over source-view training?
    \item How does GenSplat compare against diffusion-based and gaussian-based NVS methods in terms of rendering quality, policy learning, and efficiency?
    \item How does the core 3D-prior distillation component impact geometric consistency, rendering fidelity, and the success rate of downstream policies?
    \item What is the impact of increasing the number of rendered views per trajectory on view generalization and data efficiency?
\end{enumerate}

\begin{figure*}[t]
    \centering
    \includegraphics[width=\linewidth]{fig/fig_real.pdf}
    \caption{\textbf{Overview of the experiment setup and tasks.} We design six manipulation tasks for real-world evaluation. 
    }
    \label{fig:real}
    \vspace{-2ex}
\end{figure*}


\noindent \textbf{Experimental Setup.} 
Our hardware suite (\cref{fig:real}) comprises a 7-DoF Franka Research 3 arm equipped with a Robotiq 2F-85 gripper. Visual observations are captured by three standard RGB cameras: two providing external views and one mounted on the robot's wrist. 
We evaluate our framework across six manipulation tasks. For each task, we collect 100 expert demonstrations via teleoperation. To ensure rigorous evaluation, the dataset incorporates rich spatial and visual diversity through systematic variations in object positions and colors.

\subsection{RQ1: Robustness to Camera Perturbations}
\label{sec:Robustness to Camera Perturbations}
This experiment quantitatively investigates whether densifying the spatial support via GenSplat systematically improves policy robustness against camera perturbations of varying magnitudes.

\noindent \textbf{Viewpoint re-rendering strategy.}
We employ GenSplat to reconstruct the scene and synthesize novel-view trajectories while strictly preserving temporal alignment and kinematic action labels. For each demonstration, we define the pivot point as the midpoint of the shortest line segment connecting the optical axes of the two external cameras. 
Following~\cite{tian2024view,shi2025nvspolicy}, novel views are generated by rotating the camera about a designated world axis ($\mathcal{X}$ or $\mathcal{Y}$) by $\theta \in [-30^\circ, 30^\circ]$, followed by a forward translation $\mu \in [0.5, 6]$ cm along the camera's $+\mathcal{Z}$ axis. This translation crucially prevents severe rotations from pushing the viewpoint outside the densified spatial support (\textit{i.e.}, beyond valid splat coverage). For each original demonstration, we sample exactly one unique perturbation pair $(\theta,\mu)$ to generate a single augmented trajectory.


\noindent \textbf{Perturbation analysis protocol.}
% We evaluate each policy under three controlled perturbation regimes applied to the camera pose during deployment:
% (1) \textbf{Small} perturbations: Rotation $\theta_{\max}\in[3^\circ,6^\circ]$ and forward translation $\mu\in[0.5,2]\text{ cm}$.
% (2) \textbf{Medium} perturbations: Rotation $\theta_{\max}\in[8^\circ,15^\circ]$ and translation $\mu\in[2,4]\text{ cm}$.
% (3) \textbf{Large} perturbations: Rotation $\theta_{\max}\in[18^\circ,30^\circ]$ and translation $\mu\in[4,6]\text{ cm}$.
% We compare two policy variants: the baseline trained solely on source data, and with GenSplat-augmented novel views data. The augmentation uses GenSplat's feed-forward Gaussian splatting to synthesize geometrically consistent multi-view trajectories. Both policies are evaluated via real robot rollouts across all three perturbation levels.
We deploy the trained policies under three controlled spatial perturbation regimes during physical rollouts: 
(1) \textbf{Small}: Rotation $\theta_{\max}\in[3^\circ,6^\circ]$ and translation $\mu\in[0.5,2]$ cm. 
(2) \textbf{Medium}: Rotation $\theta_{\max}\in[8^\circ,15^\circ]$ and translation $\mu\in[2,4]$ cm. 
(3) \textbf{Large}: Rotation $\theta_{\max}\in[18^\circ,30^\circ]$ and translation $\mu\in[4,6]$ cm. We compare baselines trained strictly on source data against our policies trained on the expanded manifold.

As visualized in~\cref{fig:Q2}, GenSplat augmentation yields consistent and substantial improvements in policy robustness across all six manipulation tasks and all perturbation levels, mitigating the sharp degradation typically caused by out-of-distribution camera viewpoints. Aggregated across all tasks, GenSplat yields substantial relative performance improvements across all perturbation levels for \(\pi_0\): +\textbf{6.7}\% for small, +\textbf{8.9}\% for medium, and +\textbf{13.3}\% for large perturbations. The benefit is most pronounced for DP under large perturbations, where source-view-only policies suffer significant collapse, and GenSplat achieves a remarkable \textbf{56.0}\% relative gain (from 27.78\% to \textbf{43.33}\%). Crucially, the augmented policies maintain parity with the baselines in the unperturbed (source) setting, confirming that our geometrically consistent novel views do not compromise the original data distribution. These results establish GenSplat as an effective, architecture-agnostic strategy for training viewpoint generalizable policies.

\begin{figure*}[t]
    \centering
    \includegraphics[width=\linewidth]{fig/fig_Q2.pdf}
        \caption{\textbf{Policy robustness to camera pose perturbations.} We evaluate 30 episodes per task on 6 real-world tasks under increasing camera perturbations and report mean success rates (in \%). Policies trained with GenSplat augmented novel views data consistently outperform those trained only on source views for both \(\pi_{0}\) and Diffusion Policy. 
        % Quantitative results are detailed in Appendix \textcolor{red}{B.1}.
    }
    % 匿名版是 -4ex
    \vspace{-2ex}
    \label{fig:Q2}
\end{figure*}

\vspace{-2ex}
\subsection{RQ2: Comprehensive Comparison with SOTA}
We address RQ2 by evaluating GenSplat across three critical axes: (1) rendering quality, (2) downstream policy performance, and (3) synthesis efficiency. We compare GenSplat against leading NVS methods including ZeroNVS~\cite{sargent2024zeronvs}, VISTA~\cite{tian2024view}, SEVA~\cite{zhou2025stable}, InstantSplat~\cite{fan2024instantsplat}, NoPoSplat~\cite{ye2024no} and AnySplat~\cite{jiang2025anysplat}.

\noindent \textbf{Geometric Consistency \& Rendering Quality.} 
We first validate the underlying 3D structural fidelity. As illustrated in~\cref{fig:2.1}, GenSplat accurately infers dense depth and normal maps. While the distilled pseudo-labels inherently suffer from over-smoothing, our unified photometric and geometric optimization effectively recovers sharp object boundaries and fine-grained topologies. By leveraging RGB photometric supervision with 3D-priors strictly acting as structural regularizers, GenSplat successfully reconstructs high-frequency details.


\begin{figure}[t]
    \centering
    \includegraphics[width=1.0\linewidth]{fig/RQ2.1.pdf}
    \caption{\textbf{Qualitative 3D geometry reconstruction on the DROID dataset.} We visualize pairs of reference targets (ground truth or pseudo-labels) against GenSplat's predicted RGB, depth, and normal maps. GenSplat accurately recovers structural high-frequency details and maintains strict spatial alignment without explicit camera calibration. 
    Extended visual comparisons are provided in Appendix \textcolor{red}{B}.}
    \vspace{-3ex}
    \label{fig:2.1}
\end{figure}

\begin{figure}[t]
    \centering
    \includegraphics[width=1.0\linewidth]{fig/RQ2.2.pdf}
    \caption{\textbf{Qualitative comparison of NVS.} 
    We compared GenSplat with diffusion- and Gaussian-based methods under the same input example. GenSplat shows a clear advantage in 3D geometric consistency modeling while effectively eliminating floating artifacts and visual artifacts. Note: VISTA is built on ZeroNVS and is fine-tuned on a 3k trajectory subset of DROID~\cite{khazatsky2024droid}, yielding significant improvements in NVS.
    }
    \vspace{-4ex}
    \label{fig:2.2}
\end{figure}

Building upon this robust geometry, GenSplat demonstrates a significant advantage in synthesizing high-fidelity novel views. As depicted in~\cref{fig:2.2}, diffusion-based baselines frequently violate strict epipolar geometry, resulting in severe spatial hallucinations and metric distortions (\textit{e.g.}, the bent faucet). Conversely, conventional Gaussian-based methods lacking prior distillation are plagued by severe floater artifacts and boundary tearing~\cite{ye2024no}, fundamentally degrading visual quality. In stark contrast, GenSplat renders scenes with photorealism and absolute structural integrity. This geometrically consistent, high-fidelity synthesis is paramount, providing the clean, low-variance multi-view data necessary for reliable downstream policy learning.

% As depicted in~\cref{fig:2.2}, diffusion-based baselines frequently fail to maintain strict multi-view consistency, resulting in undesirable geometric distortions (\textit{e.g.}, the bent faucet) and hallucinated spatial arrangements. Conversely, conventional Gaussian-based methods without prior distillation are plagued by severe floater artifacts and boundary tearing~\cite{fan2024instantsplat}, which fundamentally degrade visual quality. In stark contrast, GenSplat renders scenes with photorealism and absolute structural integrity, effectively preserving object geometry while eliminating visual noise. This geometrically consistent, high-fidelity rendering is crucial, as it provides the clean and reliable multi-view data necessary for robust downstream policy learning. 


\definecolor{lightgreen}{HTML}{EDF8FE}
% \renewcommand{\arraystretch}{1.1}
\definecolor{grey}{RGB}{235,235,235}

\begin{table*}[t]
    \centering
    \resizebox{\linewidth}{!}{
        \setlength{\tabcolsep}{0.5em}
        \begin{tabular}{l | c | c c c c c c | c}
        \toprule
        \rowcolor{lightgreen}
        \multicolumn{1}{l|}{Method / Task} & 
        \multicolumn{1}{c|}{\begin{tabular}[c]{@{}c@{}}Avg.\\ Success↑\end{tabular}} & 
        \multicolumn{1}{c}{\begin{tabular}[c]{@{}c@{}}Pick up\\ Banana\end{tabular}} & 
        \multicolumn{1}{c}{\begin{tabular}[c]{@{}c@{}}Stack\\ Block\end{tabular}} & 
        \multicolumn{1}{c}{\begin{tabular}[c]{@{}c@{}}Insert\\ Ring\end{tabular}} & 
        \multicolumn{1}{c}{\begin{tabular}[c]{@{}c@{}}Place \\Can \end{tabular}} & 
        \multicolumn{1}{c}{\begin{tabular}[c]{@{}c@{}}Put Blocks\\ into Plate\end{tabular}} & 
        \multicolumn{1}{c|}{\begin{tabular}[c]{@{}c@{}}Stack\\ Cups\end{tabular}} & 
        \multicolumn{1}{c}{\begin{tabular}[c]{@{}c@{}}Synthesis\\ Speed\end{tabular}} \\
        \midrule

        \rowcolor{grey}
        Baseline & 48.9 & 73.3 & 60.0 & 33.3 & 26.7 & 53.3 & 46.7 & - \\

        VISTA~\cite{tian2024view} & 37.8 & 76.7 & 43.3 & 36.7 & 6.7 & 40.0 & 23.3 & 2.81 s \\

        SEVA~\cite{zhou2025stable} & 57.7 & 83.3 & 70.0 & 43.3 & 30.0 & 63.3 & 56.7 & 50.00 s \\

        InstantSplat~\cite{fan2024instantsplat} & 55.0 & 80.0 & 66.7 & 40.0 & 33.3 & 60.0 & 50.0 & 85.20 s \\

        NoPoSplat~\cite{ye2024no} & 36.7 & 73.3 & 40.0 & 33.3 & 10.0 & 36.7 & 26.7 & \textbf{0.05} s \\

        AnySplat~\cite{jiang2025anysplat} & 56.1 & 80.0 & 66.7 & 43.3 & 30.0 & 63.3 & 53.3 & 0.16 s \\

        \midrule 
        GenSplat(ours) & \textbf{62.2} & \textbf{86.7} & \textbf{73.3} & \textbf{46.7} & \textbf{40.0} & \textbf{66.7} & \textbf{60.0} & 0.14 s \\ 
        \bottomrule 

        \end{tabular}
    }
    \caption{\textbf{Policy Performance and Synthesis Efficiency.} We compare various NVS methods by training \(\pi_0\) on augmented datasets and report mean success rate and synthesis speed. Each task was evaluated 30 episodes under the \textbf{Large} perturbation level. Synthesis speed evaluated on the same computing device.}
    \label{tab:RQ4}
    \vspace{-6ex}
\end{table*}

\noindent \textbf{Downstream Policy Performance.} 
As detailed in~\cref{tab:RQ4}, GenSplat achieves the highest average success rate of \textbf{62.2}\% under large camera perturbations, significantly outperforming all competing NVS paradigms. Crucially, the results reveal that low-quality data augmentation actively harms policy learning: both the diffusion-based VISTA~\cite{tian2024view} (37.8\%) and the standard feed-forward NoPoSplat~\cite{ye2024no} (36.7\%) severely degrade performance below the unaugmented Baseline (48.9\%). Their pronounced geometric distortions and rendering artifacts act as noise, contaminating the visuomotor mapping (see~\cref{fig:2.2}). While SEVA~\cite{zhou2025stable} (57.7\%) and InstantSplat~\cite{fan2024instantsplat} (55.0\%) offer improvements, their sparse-view reconstructions still exhibit localized holes and inconsistencies that bottleneck execution accuracy. Compared to the closest feed-forward baseline, AnySplat~\cite{jiang2025anysplat} (56.1\%), GenSplat provides a substantial absolute gain of \textbf{+6.1}\%. This confirms that our 3D-prior distillation explicitly enforces the spatial smoothness necessary to generate physically reliable, high-fidelity training signals.

\noindent \textbf{Synthesis Efficiency.} 
% For scalable data augmentation in robotics, computational efficiency is crucial. GenSplat operates at a highly efficient 0.14 seconds per frame. It completely bypasses the prohibitive computational bottlenecks of iterative diffusion sampling and per-scene optimization. While NoPoSplat~\cite{ye2024no} is marginally faster (0.05 s), this speed comes at the unacceptable cost of catastrophic policy degradation. 
For scalable data augmentation in robotics, computational efficiency is a hard constraint. GenSplat operates at a highly efficient 0.14 seconds per frame, completely bypassing the prohibitive computational bottlenecks of iterative diffusion sampling and per-scene optimization. While NoPoSplat~\cite{ye2024no} is marginally faster, this speed incurs the unacceptable cost of catastrophic policy degradation. 
Ultimately, GenSplat achieves remarkable generation speed without sacrificing geometric consistency, positioning our framework as a highly scalable engine for learning view-generalized robotic policies.


\vspace{-2ex}
\subsection{RQ3: Ablation of 3D-Prior Distillation}
To address RQ3, we conduct an ablation study isolating the contributions of the permutation-equivariant (P.E.) architecture \cite{wang2025pi} and the core 3D-prior distillation module. \cref{tab:ablation} details the quantitative impact of these components on both novel view synthesis fidelity and downstream policy robustness under challenging viewpoint shifts.
% \cref{tab:ablation} quantifies their impact on both novel-view synthesis fidelity and downstream policy robustness under severe spatial perturbations.
Replacing the baseline VGGT with the P.E. architecture improves reconstruction. By mitigating order-dependent artifacts from uncalibrated inputs, it boosts policy success from 66.0\% to 70.0\% under large perturbations, and 30.0\% to 40.0\% under extreme shifts. 
Furthermore, integrating 3D-prior distillation provides essential geometric regularization for pose-free settings, achieving the best PSNR and LPIPS. This strict structural consistency translates directly to physical robustness, maximizing policy success to \textbf{74.0\%} (large) and \textbf{46.0\%} (extreme). By explicitly suppressing floating artifacts, the 3D prior prevents geometric collapse, delivering the high-fidelity manifold data necessary to induce view-generalized behaviors.
\renewcommand{\arraystretch}{1.0} % 设置行距为1.5倍

\definecolor{darkgreen}{HTML}{25C445}
\definecolor{darkred}{HTML}{DC143C}
\definecolor{lightgrey}{HTML}{E8E8E8}
\definecolor{final}{HTML}{62EF7E}
\definecolor{base}{HTML}{BDBDBD}

\begin{table}[t]
    \centering
    \resizebox{0.9\linewidth}{!}{
        \setlength{\tabcolsep}{1.0em} 
        \begin{tabular}{ccc|cccc} % 修改此处：增加两列，变为 ccc|cccc
        \toprule
        \rowcolor{lightgreen}
        VGGT & P.E.~\cite{wang2025pi} & \begin{tabular}[c]{@{}c@{}} 3D-Prior \end{tabular} & \begin{tabular}[c]{@{}c@{}} PSNR↑ \end{tabular} & \begin{tabular}[c]{@{}c@{}}LPIPS↓\end{tabular} & \begin{tabular}[c]{@{}c@{}} $30^{\circ}$, 6cm \end{tabular} & \begin{tabular}[c]{@{}c@{}} $60^{\circ}$, 12cm \end{tabular} \\ 
        \midrule
         
        {\color{darkgreen} \Checkmark} & {\color{darkred} \XSolidBrush} & {\color{darkred} \XSolidBrush} & 22.35 & 6.01 & 66.0\% & 30.0\% \\
         
        {\color{darkred} \XSolidBrush} & {\color{darkgreen} \Checkmark} & {\color{darkred} \XSolidBrush} & 25.87 & 5.03 & 70.0\% & 40.0\% \\
        
        {\color{darkred} \XSolidBrush} & {\color{darkgreen} \Checkmark} & {\color{darkgreen} \Checkmark} & \textbf{26.53} & \textbf{3.72} & \textbf{74.0\%} & \textbf{46.0\%} \\
         
         \bottomrule
        \end{tabular}
    }
    \caption{\textbf{Ablation of GenSplat components.} We evaluate the impact of the permutation-equivariant (P.E.) architecture and 3D-prior distillation on novel view rendering quality (PSNR, LPIPS) and policy average success rates under \textbf{Large} ($30^\circ$, 6cm) and \textbf{Extreme} ($60^\circ$, 12cm) camera perturbations.}
    \label{tab:ablation}
    \vspace{-7ex}
\end{table}


\vspace{-4.5ex}
\begin{figure}[h]
    \centering
    \begin{minipage}[c]{0.48\linewidth}
    As illustrated in~\cref{fig:anysplat_cmp}, we utilize 3D point maps as structural pseudo-labels, unlike AnySplat~\cite{jiang2025anysplat}, which relies on scalar depth maps. Depth-based supervision often fails to maintain rigorous multi-view consistency, rendering the model incapable of reconstructing flat regions and leading to severe structural collapse of the desk. By directly encoding explicit 3D coordinates, point maps offer higher precision and enforce spatial smoothness~\cite{shi2025revisiting,wang2024dust3r}. This effectively prevents geometric collapse, recovering a remarkably accurate and flat desk surface. More visual examples are in Appendix \textcolor{red}{B.2}.
    \end{minipage}
    \hfill
    % \vspace{-2ex}
    \begin{minipage}[c]{0.48\linewidth}
        \includegraphics[width=\linewidth]{fig/fig_3.1.pdf}
        \caption{\textbf{Qualitative scene reconstruction comparison.} AnySplat (left) exhibits geometric irregularities on flat surfaces. GenSplat (right) reconstructs accurate, smooth 3D geometry.}
        \label{fig:anysplat_cmp}
    \end{minipage}
    \vspace{-4ex}
\end{figure}


%%%%%%%%%%%%%%%%%%%%%%%%%%%%%%%%%%%%%%%%%%%%%%%%%%%%%%%%%%%%%%%%%%%%%%%%%%%%%%%%%%%%%%%%%%%%
% 在匿名版本中是要加上的！
% \vspace{-3ex}
\subsection{RQ4: Impact of the Number of Rendered Views}
For RQ4, we quantify the trade-off between the number of synthesized novel views and policy generalization. Fixing the source trajectories at 100, we systematically vary the view augmentation factor \(N \in \{1, 2, 3, 4\}\) and evaluate the resulting Diffusion Policies~\cite{chi2023diffusion}. According to~\cref{fig:Q3}, under a \textbf{Medium} perturbation regime, increasing the number of rendered views consistently improves policy success rates. Notably, the most substantial performance leap occurs immediately upon introducing the first augmented view ($N=1$).

% \vspace{-2ex}
\begin{figure*}[h!]
    \centering
    \vspace{-3ex}
    \includegraphics[width=0.95\linewidth]{fig/fig_Q3.pdf}
    \vspace{-3ex}
    \caption{\textbf{Effect of Rendered Views on Viewpoint Robustness.}
    % Diffusion Policy evaluated across six tasks under \textbf{Medium} perturbations. 
    We report the mean success rate and standard deviation across three independent trials, with 30 execution episodes per task in each run. 
    % The baseline (100S) uses 100 source trajectories; 100S+$k$G indicates augmentation with $k \in \{100, 200, 300, 400\}$ GenSplat-rendered views, corresponding to 1 to 4 novel views per source trajectory.
    }
    \vspace{-3ex}
    \label{fig:Q3}
\end{figure*}


% \vspace{-4ex}
\begin{table}[h]
    \centering
    % \renewcommand{\arraystretch}{1.1} 
    \begin{tabularx}{0.9\linewidth}{ l *{8}{>{\centering\arraybackslash}X} }
        \toprule
        \rowcolor{lightblue}
        Task & \multicolumn{4}{c}{Stack Block} & \multicolumn{4}{c}{Stack Cups} \\
        \rowcolor{lightblue}
        Perturbation & None & Small & Medium & Large & None & Small & Medium & Large \\
        \midrule
        
        \rowcolor{grey}
        100S (Baseline) & 68.0 & 58.0 & 42.0 & 30.0 & 60.0 & 50.0 & 40.0 & 26.0 \\
        100S+100G       & 70.0 & 66.0 & 62.0 & 48.0 & 62.0 & 58.0 & 50.0 & 44.0 \\
        100S+200G       & 68.0 & 70.0 & 68.0 & 58.0 & 60.0 & 62.0 & 58.0 & 50.0 \\
        100S+300G       & \textbf{72.0} & \textbf{74.0} & \textbf{74.0} & \textbf{64.0} & \textbf{64.0} & 64.0 & \textbf{64.0} & \textbf{56.0} \\
        100S+400G       & 70.0 & 74.0 & 72.0 & 62.0 & 62.0 & \textbf{66.0} & 64.0 & 54.0 \\
        
        \bottomrule
    \end{tabularx}
    \caption{\textbf{Cross-Perturbation Analysis.} We report the success rates for Diffusion Policy with varying GenSplat view augmentations, evaluated on two tasks across perturbation levels from None to Large, based on 50 episodes per task.}
    \label{tab:RQ3-2}
    \vspace{-5ex}
\end{table}

% 匿名版本中是[h]
\begin{figure*}[ht!]
    \centering
    \includegraphics[width=0.9\linewidth]{fig/fig_manifold.pdf}
    \caption{
    \textbf{Densification of camera extrinsics on the observational manifold.}
    % \textbf{Geometric visualization of manifold densification.}
    }
    \vspace{-4ex}
    \label{fig:fig_manifold}
\end{figure*}


% \noindent \textbf{Analysis of Data Efficiency.}
% % Critically, the performance gains exhibit rapid saturation. \cref{tab:RQ3-2} reports the extended view augmentation analysis across all perturbation magnitudes (None to Large) on the \texttt{Stack Block} and \texttt{Stack Cups} tasks. The results confirm a consistent scaling pattern invariant to perturbation severity: after the sharp initial gain at $N=1$, the marginal benefits steadily diminish, saturating around $N=3$, with $N=4$ yielding negligible improvements.
% % This cross-validation affirms that a moderate set of geometrically consistent GenSplat augmentations ($N=3$) is sufficient to unlock strong view generalization, effectively avoiding the computational overhead of excessive rendering.
\noindent \textbf{Analysis of Data Efficiency.}
Critically, performance gains saturate rapidly. \cref{tab:RQ3-2} extends this analysis across all perturbation severities on the \texttt{Stack Block} and \texttt{Stack Cups} tasks, revealing a consistent scaling law: after a sharp initial leap at $N=1$, marginal benefits diminish and saturate near $N=3$. Since $N=4$ yields negligible improvements, a moderate augmentation density ($N=3$) is proven sufficient to unlock robust view generalization, effectively circumventing the overhead of excessive rendering.

\noindent \textbf{Geometric Interpretation.}
To geometrically explain this saturation, \cref{fig:fig_manifold} visualizes a single camera's extrinsics across 600 real-world demonstrations. The initial augmentation ($N=1$) sharply expands the sparse source distribution ($N=0$), driving the primary performance leap. Subsequent views ($N \in \{2, 4\}$) merely densify this expanded spherical manifold without enlarging its boundary. Once topological continuity is achieved (near $N=3$), the policy interpolates robustly, rendering further augmentations computationally redundant.
